% Options for packages loaded elsewhere
% Options for packages loaded elsewhere
\PassOptionsToPackage{unicode}{hyperref}
\PassOptionsToPackage{hyphens}{url}
\PassOptionsToPackage{dvipsnames,svgnames,x11names}{xcolor}
%
\documentclass[
  letterpaper,
  DIV=11,
  numbers=noendperiod]{scrartcl}
\usepackage{xcolor}
\usepackage{amsmath,amssymb}
\setcounter{secnumdepth}{-\maxdimen} % remove section numbering
\usepackage{iftex}
\ifPDFTeX
  \usepackage[T1]{fontenc}
  \usepackage[utf8]{inputenc}
  \usepackage{textcomp} % provide euro and other symbols
\else % if luatex or xetex
  \usepackage{unicode-math} % this also loads fontspec
  \defaultfontfeatures{Scale=MatchLowercase}
  \defaultfontfeatures[\rmfamily]{Ligatures=TeX,Scale=1}
\fi
\usepackage{lmodern}
\ifPDFTeX\else
  % xetex/luatex font selection
\fi
% Use upquote if available, for straight quotes in verbatim environments
\IfFileExists{upquote.sty}{\usepackage{upquote}}{}
\IfFileExists{microtype.sty}{% use microtype if available
  \usepackage[]{microtype}
  \UseMicrotypeSet[protrusion]{basicmath} % disable protrusion for tt fonts
}{}
\makeatletter
\@ifundefined{KOMAClassName}{% if non-KOMA class
  \IfFileExists{parskip.sty}{%
    \usepackage{parskip}
  }{% else
    \setlength{\parindent}{0pt}
    \setlength{\parskip}{6pt plus 2pt minus 1pt}}
}{% if KOMA class
  \KOMAoptions{parskip=half}}
\makeatother
% Make \paragraph and \subparagraph free-standing
\makeatletter
\ifx\paragraph\undefined\else
  \let\oldparagraph\paragraph
  \renewcommand{\paragraph}{
    \@ifstar
      \xxxParagraphStar
      \xxxParagraphNoStar
  }
  \newcommand{\xxxParagraphStar}[1]{\oldparagraph*{#1}\mbox{}}
  \newcommand{\xxxParagraphNoStar}[1]{\oldparagraph{#1}\mbox{}}
\fi
\ifx\subparagraph\undefined\else
  \let\oldsubparagraph\subparagraph
  \renewcommand{\subparagraph}{
    \@ifstar
      \xxxSubParagraphStar
      \xxxSubParagraphNoStar
  }
  \newcommand{\xxxSubParagraphStar}[1]{\oldsubparagraph*{#1}\mbox{}}
  \newcommand{\xxxSubParagraphNoStar}[1]{\oldsubparagraph{#1}\mbox{}}
\fi
\makeatother


\usepackage{longtable,booktabs,array}
\newcounter{none} % for unnumbered tables
\usepackage{calc} % for calculating minipage widths
% Correct order of tables after \paragraph or \subparagraph
\usepackage{etoolbox}
\makeatletter
\patchcmd\longtable{\par}{\if@noskipsec\mbox{}\fi\par}{}{}
\makeatother
% Allow footnotes in longtable head/foot
\IfFileExists{footnotehyper.sty}{\usepackage{footnotehyper}}{\usepackage{footnote}}
\makesavenoteenv{longtable}
\usepackage{graphicx}
\makeatletter
\newsavebox\pandoc@box
\newcommand*\pandocbounded[1]{% scales image to fit in text height/width
  \sbox\pandoc@box{#1}%
  \Gscale@div\@tempa{\textheight}{\dimexpr\ht\pandoc@box+\dp\pandoc@box\relax}%
  \Gscale@div\@tempb{\linewidth}{\wd\pandoc@box}%
  \ifdim\@tempb\p@<\@tempa\p@\let\@tempa\@tempb\fi% select the smaller of both
  \ifdim\@tempa\p@<\p@\scalebox{\@tempa}{\usebox\pandoc@box}%
  \else\usebox{\pandoc@box}%
  \fi%
}
% Set default figure placement to htbp
\def\fps@figure{htbp}
\makeatother


% definitions for citeproc citations
\NewDocumentCommand\citeproctext{}{}
\NewDocumentCommand\citeproc{mm}{%
  \begingroup\def\citeproctext{#2}\cite{#1}\endgroup}
\makeatletter
 % allow citations to break across lines
 \let\@cite@ofmt\@firstofone
 % avoid brackets around text for \cite:
 \def\@biblabel#1{}
 \def\@cite#1#2{{#1\if@tempswa , #2\fi}}
\makeatother
\newlength{\cslhangindent}
\setlength{\cslhangindent}{1.5em}
\newlength{\csllabelwidth}
\setlength{\csllabelwidth}{3em}
\newenvironment{CSLReferences}[2] % #1 hanging-indent, #2 entry-spacing
 {\begin{list}{}{%
  \setlength{\itemindent}{0pt}
  \setlength{\leftmargin}{0pt}
  \setlength{\parsep}{0pt}
  % turn on hanging indent if param 1 is 1
  \ifodd #1
   \setlength{\leftmargin}{\cslhangindent}
   \setlength{\itemindent}{-1\cslhangindent}
  \fi
  % set entry spacing
  \setlength{\itemsep}{#2\baselineskip}}}
 {\end{list}}
\usepackage{calc}
\newcommand{\CSLBlock}[1]{\hfill\break\parbox[t]{\linewidth}{\strut\ignorespaces#1\strut}}
\newcommand{\CSLLeftMargin}[1]{\parbox[t]{\csllabelwidth}{\strut#1\strut}}
\newcommand{\CSLRightInline}[1]{\parbox[t]{\linewidth - \csllabelwidth}{\strut#1\strut}}
\newcommand{\CSLIndent}[1]{\hspace{\cslhangindent}#1}



\setlength{\emergencystretch}{3em} % prevent overfull lines

\providecommand{\tightlist}{%
  \setlength{\itemsep}{0pt}\setlength{\parskip}{0pt}}



 


\KOMAoption{captions}{tableheading}
\makeatletter
\@ifpackageloaded{caption}{}{\usepackage{caption}}
\AtBeginDocument{%
\ifdefined\contentsname
  \renewcommand*\contentsname{Table of contents}
\else
  \newcommand\contentsname{Table of contents}
\fi
\ifdefined\listfigurename
  \renewcommand*\listfigurename{List of Figures}
\else
  \newcommand\listfigurename{List of Figures}
\fi
\ifdefined\listtablename
  \renewcommand*\listtablename{List of Tables}
\else
  \newcommand\listtablename{List of Tables}
\fi
\ifdefined\figurename
  \renewcommand*\figurename{Figure}
\else
  \newcommand\figurename{Figure}
\fi
\ifdefined\tablename
  \renewcommand*\tablename{Table}
\else
  \newcommand\tablename{Table}
\fi
}
\@ifpackageloaded{float}{}{\usepackage{float}}
\floatstyle{ruled}
\@ifundefined{c@chapter}{\newfloat{codelisting}{h}{lop}}{\newfloat{codelisting}{h}{lop}[chapter]}
\floatname{codelisting}{Listing}
\newcommand*\listoflistings{\listof{codelisting}{List of Listings}}
\makeatother
\makeatletter
\makeatother
\makeatletter
\@ifpackageloaded{caption}{}{\usepackage{caption}}
\@ifpackageloaded{subcaption}{}{\usepackage{subcaption}}
\makeatother
\usepackage{bookmark}
\IfFileExists{xurl.sty}{\usepackage{xurl}}{} % add URL line breaks if available
\urlstyle{same}
\hypersetup{
  pdftitle={Summaries},
  pdfauthor={Robert Esposito},
  colorlinks=true,
  linkcolor={blue},
  filecolor={Maroon},
  citecolor={Blue},
  urlcolor={Blue},
  pdfcreator={LaTeX via pandoc}}


\title{Summaries}
\author{Robert Esposito}
\date{}
\begin{document}
\maketitle


\section{Article Summaries}\label{article-summaries}

Note: If I can't find the article yet, I just put the abstract.

\subsection{Stephens (1989)}\label{stephens1989language}

Portuguese in USA. Arrivals began in 1820s, picked up in 1830s with
arrival fo Lusitanian and Azorean. Congregated to New England,
specifically MA. 1850s, to California (San Fran).

Late 19th/early 20th century, big influx to CA, MA, NY, \& NJ.
Mid-1970s, major immigration event.

Early immigrants = high illiteracy, low English.

Azorean Portuguese is stigamized (Middleton, 1986).

Began arriving to MNeark in large groups in 1920s (Cunha, 1981).

Bilingual programs (I guess school programs?) \& Sunday schools. 15
private Portuguese schools, similar to Hebrew schools.

\subsection{\texorpdfstring{(\textbf{caravalho201014?})}{(caravalho201014?)}}\label{caravalho201014}

Portuguese immigrants mainly from Azores Islands, mainland Portugual,
and Cape Verdean, smaller population from Mozambique and Angola. High
concentrations in MA. Over last 20 years, Cape Verdeans \& Brazilians
account for most PT speakers.

14,829 students taking Portuguese classes in K-12 schools in 2008
(Vicente \& Pimenta, 2008).

Modern Language Association in 2006 reported 22.4\% increase of
availability of Portuguese instruction between 2002 and 2006.

US Census underestimates PT speakers in USA. Ministry of Foreign Affairs
in Brazil report much higher estimate. Due to undocumented immigrants.

National Bilingual Education Act of 1968 resulted in PT-EN bilingual
education initiatives in 1970s.

Republic of Cape Verde immigrants speak a Portuguese-based creole and
Portuguese. They typically are more receptive bilinguals than e.g.,
Azorean or Madeiran speakers in the US (Ferreira 2005; need to find
bibtex). Cape Verdean-Americans tend to be L1 Crole L1 EN/L2 EN, and
then may learn PT as L3 in school.

Brazilian immigration in 1984-1987, again in 1990s. Young, educated,
middle-class men and women on tourist visas and overstayed (Bensabat-Ott
2000; need to find bibtex). Later, working class immigrants from small
towns and rural areas who entered via US-Mexico border.

Luso-American Foundation (FLAD, Fundacao Luso-Americana de
Desenvolvimento) identified
\texttt{r\ 52\ +\ 25\ +\ 20\ +\ 17\ +\ 13\ +\ 10\ +\ 5\ +1} K-12 public
and community schools that offer Portuguese. 20 in NJ, 10 in NY, 13 in
CT; weekend Portuguese language schools in NJ (Vicente \& Pimenta,
2008).

Dual-immersion EN-PT program in K-8 school in MA resulted in greater
awareness of dialectal differences among teachers, parents, and staff
(Rubinstein-Avila, 2002).

Dialectal differences do not result in communicative breakdown between
students (Rubinstein-Avila, 2002; Ferreira, 2005).

\subsection{Osborne (2021)}\label{osborne2021perception}

Mid vowel contrasts as the last to be acquired by L1 speakers
(Matzenauer \& Costa, 2017). Low functional load with few minimal pairs
(Alves, 2008).

12 BR-PT heritage speakers. Completed BLP (all except one was
English-dominant, but pretty large range).

AXB discrimination. Check article for list of stimuli.

2AFC identification task.

Heritage speakers accurately categorize PT mid vowels, but higher EN
dominance predicted more trouble with /ɔ/-/o/. Amengual (2015) found the
same asymmetry for ES-Catalan bilinguals.

\subsection{Rothman \& Judy (2014)}\label{rothman2014portuguese}

No studies examining Cape Verdean Portuguese as a heritage language.
(Although, as noted by (\textbf{caravalho201014?}), they would be
learning Cape Verdean Creole at home).

PT-PT stops are pronounced as fricatives; BR-PT pronounces as plosives.
Would be interesting to see if PT-PT maintains fricatives in Newark
given input from BR-PT and English.

Summary of heritage speakers learning in college, focused on dialectal
differences (Silva, 2010).

\subsection{Silva (2010)}\label{silva2010starting}

Developing a heritage speaker track in college for PT. Focus on
dialectal differences.

\subsection{Galves et al. (2005)}\label{galves2005syntax}

Syntactic position of clitic pronouns (e.g., \emph{me, te, se, o, lhe})
in EP and BP arise from both syntactic and morphophonological factors.
The authors argue that EP clitics attach to Infl (inflectional category)
and BP clitics attach to the verb (V) in the syntactic structure.
Additionally, the authors argue that EP does not permit clitics in
sentence-initial position, whereas BP does not have this restriction.
This results in a strong preference for enclisis (post-verbal clitic) in
EP (\emph{chamou-me}), and a preference for proclisis (pre-verbal
clitic) in BP (\emph{me chamou}). In multi-verb constructions, EP
clitics attach to the finite verb (\emph{tinham-se entendido
perfeitamente}), while in BP they tend to attach to the non-finite verb
(\emph{tinham se entendido perfeitamente})

\subsection{Amadeo (2025)}\label{amadeo2025perception}

Escudero et al.~(2009; ``A cross-dialect acoustic description of vowels:
Brazilian and European Portuguese) and Rato \& Carlet (2020) show that
BR-PT and PT-PT share some similarities in their vowel inventory.

Diaz-Granado (2011) found that /e \textasciitilde{} ɛ/ in English \&
Portuguese are not similar, so L1 En L2 Pt are expected to collapse PT
/e \textasciitilde{} ɛ/ into a single category. Tested L1 En L2 Pt and
L1 En L2 Es L3 Pt. L3 learners performed better. But there were only 12
in the first group, 5 in the second.

Kendall (2004) tested 20 beginning and advanced L1 En L2 Pt learners on
BP mid vowels. They did poorly.

Vasiliev (2013) tested L1 English L2 Pt stressed oral vowels. /e
\textasciitilde{} ɛ/ was second-easiest.

Elvin et al.~(2018) L1 En monolinguals, En-Es bilinguals tested on
discrimination task for Pt-Pt and Pt-Br. /o \textasciitilde{} ɔ/ easier
in EP than BP, which was second-easiest to discriminate. Participants
had more trouble with /i \textasciitilde{} e/ and /u \textasciitilde{}
o/.

Summary: For L1 US En, /e ɛ o c/ are hardest to discriminate; /e
\textasciitilde{} ɛ, o \textasciitilde{} ɔ, e \textasciitilde{} i, o
\textasciitilde{} u/ are most challenging minimal pairs. These
challenges are accurately predicted by perception models e.g., Escudero
2009.

NEED TO FINISH THIS SUMMARY AND ALSO FIND ALL OF THESE ARTICLES.

\section{Topic Summaries}\label{topic-summaries}

\section{Misc.}\label{misc.}

Fernanda A. Ferreira has a lot of research on heritage speakers in
college classes.

\protect\phantomsection\label{refs}
\begin{CSLReferences}{1}{0}
\bibitem[\citeproctext]{ref-amadeo2025perception}
Amadeo, M. E. de S. G. (2025). \emph{The perception of portuguese vowels
by bilingual and heritage californian english learners of portuguese}
{[}Master's thesis{]}. California State University, Fresno.

\bibitem[\citeproctext]{ref-cunha1981portugueses}
Cunha, A. (1981). Os portugueses no ironbound. \emph{Luso-Americano}.

\bibitem[\citeproctext]{ref-galves2005syntax}
Galves, C., Moraes, M. A. T., \& Ribeiro, I. (2005). Syntax and
morphology in the placement of clitics in european and brazilian
portuguese. \emph{Journal of Portuguese Linguistics}, \emph{4}(2).

\bibitem[\citeproctext]{ref-middleton1986language}
Middleton, E. (1986). \emph{Language maintenance of the
portuguese-americans}.

\bibitem[\citeproctext]{ref-osborne2021perception}
Osborne, D. M. (2021). Perception of portuguese mid vowels by heritage
speakers of brazilian portuguese. \emph{Heritage Language Journal},
\emph{18}(1), 1--32.

\bibitem[\citeproctext]{ref-rothman2014portuguese}
Rothman, J., \& Judy, T. (2014). Portuguese heritage bilingualism in the
united states. In \emph{Handbook of heritage, community, and native
american languages in the united states} (pp. 132--141). Routledge.

\bibitem[\citeproctext]{ref-rubinstein2002problematizing}
Rubinstein-Avila, E. (2002). Problematizing the {``dual''} in a
dual-immersion program: A portrait. \emph{Linguistics and Education},
\emph{13}(1), 65--87.

\bibitem[\citeproctext]{ref-silva2010starting}
Silva, G. V. (2010). On starting a course sequence for heritage learners
of portuguese. \emph{Portuguese Language Journal}, \emph{4}, 1--7.

\bibitem[\citeproctext]{ref-stephens1989language}
Stephens, T. M. (1989). Language maintenance and ethnic survival: The
portuguese in new jersey. \emph{Hispania}, \emph{72}(3), 716--720.

\bibitem[\citeproctext]{ref-vicente2008promoccao}
Vicente, A. L., \& Pimenta, M. (2008). Promo{ç}{ã}o da l{ı́}ngua
portuguesa no mundo. \emph{FLAD, Lisboa}.

\end{CSLReferences}




\end{document}
